\scanpage{32}\mspage{31}

Ainsi $\mathrm{Cont}'_4$ dit que deux figures $F,\unsure{G|C}$ sont compatibles sss pour tt strate ouverte $\deleted{\unsure{A|X}}U$ de $F$ et $V$ de $G$, \deleted{telles que $U\ne V\unsure{_{\!a}|\,a}$, on ait} on ait $U=V$ ou $U$ et $V$ sont $C$-disjoints.

\uline{12. Juin} Depuis hier j'ai pas mal réfléchi à la combinatoire des pseudo-disques topologiques. Il est apparu que si $\deleted{\unsure{L|X|D}}$ $(D,\partial D)$ est un tel pseudo-disque (\unsure{en|au} un sens convenable), alors les ``lieux'' sur $D-\partial D$ ne sont \uline{pas} seulement des parties (compactes contractibles etc) de $D-\partial D$, satisfaisant à certaines conditions, mais elles doivent être accompagnées. d'une donnée supplémentaire, une \unsure{classe|donnée} d'équivalence de ``\unsure{cordes|corde}'' les joignant à $\partial D$. Du coup, cela montre que l'interprétation \deleted{\unsure{que}} que j'avais en vue, d'une structure de contrée en termes \unsure{dé|de} donnée de l'ens $\mathcal L$ des lieux, et de certains ensembles de parties de $\mathcal L$, les ``figures'', n'est pas
